\documentclass[aspectratio=169,utf8,t]{beamer}

% Idioma e Codificacao
\usepackage[brazilian]{babel}
\usepackage{booktabs}
\usepackage{array}
\usepackage{listings}
\usepackage{amsmath}
\usepackage{amssymb}
\usepackage{graphicx}
\usepackage{tikz}
\usetikzlibrary{shapes,arrows,positioning,calc}

% Carrega o Tema Modelo dos Scouts (Widescreen 16:9, Nunito Sans)
\usetheme{scouts}

% Configuracao para Codigo-Fonte (Bash / Python / Redis)
\lstset{
    basicstyle=\fontsize{6.2pt}{7.8pt}\selectfont\ttfamily,
    keywordstyle=\bfseries\color{ScoutNavy},
    stringstyle=\color{ScoutGreen},
    commentstyle=\itshape\color{gray},
    backgroundcolor=\color{black!3},
    frame=single,
    rulecolor=\color{black!15},
    breaklines=true,
    breakatwhitespace=true,
    showstringspaces=false,
    tabsize=2,
    aboveskip=3pt,
    belowskip=2pt
}

% Metadados Institucionais
\author{Prof. Me. Reginaldo Fernandes}
\institute{Instituto Federal do Ceará -- Campus Tauá \\ Tecnologia em Análise e Desenvolvimento de Sistemas}
\title{Bancos de Dados Não-Relacionais}
\subtitle{Aula 05: Redis na Prática: Estratégias de Cache (TTL), Controle de Sessão e Filas}
\date{20 de Outubro de 2026}

\begin{document}

% ------------------------------------------------------------------------------
% Capa da Apresentacao
% ------------------------------------------------------------------------------
\changecolours[inverse]{ScoutGreen}
\frame[plain]{\titlepage}

% ------------------------------------------------------------------------------
% Sumario Executivo da Aula
% ------------------------------------------------------------------------------
\changecolours{ScoutGreen}
\begin{frame}{Agenda do Encontro}
  \vspace{0.4cm}
  \tableofcontents
\end{frame}

% ------------------------------------------------------------------------------
% Modulos da Aula
% ------------------------------------------------------------------------------
\changecolours{ScoutGreen}
\section[Projeto FastCommerce]{Projeto FastCommerce e a Persistência Poliglota}

% ------------------------------------------------------------------------------
\twocolframe[title={Persistência Poliglota na Arquitetura Moderna},leftcol=7.1cm,rightcol=6.9cm]{%
  \alert{O Fim do Monolito de Dados Único}\par
  \itemseps[topsepi=0.2mm,itemsepi=0.5mm,labelsep=1.5mm,leftmargini=2.5mm]
  \begin{itemize}
    \item \textbf{Cenário Tradicional:} Um único banco relacional sofrendo com leitura pesada de catálogo e filas concorrentes.
    \item \textbf{Persistência Poliglota:} Adoção de tecnologias NoSQL distintas especializadas em cada regra de negócio.
    \item \textbf{Solução em Camadas:}
      \begin{itemize}
        \item \textbf{Redis (RAM):} Sessões rápidas, cache expirável e filas transitórias de checkout.
        \item \textbf{MongoDB (Disco/BSON):} Catálogo flexível e histórico de pedidos (Aula 06).
      \end{itemize}
  \end{itemize}
}{%
  \alert{O Projeto Transversal da Disciplina}\par
  \itemseps[topsepi=0.2mm,itemsepi=0.5mm,labelsep=1.5mm,leftmargini=2.5mm]
  \begin{itemize}
    \item \textbf{FastCommerce NoSQL:} Aplicação web real desenvolvida em sala de aula.
    \item \textbf{Ciclo Metodológico:}
      \begin{enumerate}
        \item \textbf{Projeto:} Modelagem de chaves e tipos no Redis.
        \item \textbf{Implementação:} Código Python com API e interface web.
        \item \textbf{Uso Prático:} Navegação, telemetria e filas.
      \end{enumerate}
    \item \textbf{Continuidade:} O mesmo código será estendido nas Aulas 06 a 08 rumo à AVP1.
  \end{itemize}
}

% ------------------------------------------------------------------------------
\begin{frame}{Arquitetura de Dados em Tempo Real do FastCommerce}
  \vspace{0.1cm}
  \begin{center}
    \tikz[scale=0.76, transform shape]{
      % Camada Web Front-end
      \node[draw=ScoutGreen, fill=ScoutGreen!15, thick, rounded corners, minimum width=3.2cm, minimum height=0.8cm, font=\bfseries\sffamily\scriptsize] (front) {Interface Web (HTML/JS)};

      % API Back-end
      \node[draw=ScoutNavy, fill=ScoutNavy!15, thick, rounded corners, minimum width=3.2cm, minimum height=0.8cm, font=\bfseries\sffamily\scriptsize, below=0.8cm of front] (api) {API Python (Flask / FastAPI)};

      % Redis em Memória
      \node[draw=red!80!black, fill=red!10, thick, rounded corners, minimum width=4.0cm, minimum height=1.6cm, font=\bfseries\sffamily\scriptsize, right=1.6cm of api] (redis) {
        \begin{tabular}{c}
          \textbf{REDIS (In-Memory 6379)} \\
          \tiny $\bullet$ Hashes: \texttt{cart:<session>} \\
          \tiny $\bullet$ Strings c/ TTL: \texttt{session:<id>}, \texttt{cache:catalogo} \\
          \tiny $\bullet$ Lists: \texttt{fila:pedidos}
        \end{tabular}
      };

      % Worker Background
      \node[draw=blue!80!black, fill=blue!10, thick, rounded corners, minimum width=3.4cm, minimum height=0.8cm, font=\bfseries\sffamily\scriptsize, below=1.0cm of redis] (worker) {Worker Assíncrono (\texttt{BRPOP})};

      % Conexoes
      \draw[<->, thick, ScoutGreen] (front) -- node[right, font=\tiny] {HTTP / JSON} (api);
      \draw[<->, thick, red!80!black] (api) -- node[above, font=\tiny] {Cache / Sessão / Fila} (redis);
      \draw[->, thick, blue!80!black] (redis) -- node[right, font=\tiny] {BRPOP fila:pedidos} (worker);
    }
  \end{center}
  \vspace{0.1cm}
  \begin{columns}[t]
    \begin{column}{0.48\textwidth}
      \footnotesize \textbf{Camada de Entrada (Memória):} Respostas em microssegundos ($\mu\text{s}$) para 100\% das interações ativas do usuário.
    \end{column}
    \begin{column}{0.48\textwidth}
      \footnotesize \textbf{Desacoplamento Assíncrono:} O checkout enfileira o pedido instantaneamente, liberando a interface do usuário.
    \end{column}
  \end{columns}
\end{frame}

\changecolours{ScoutGreen}
\section[Cache-Aside e TTL]{Padrão Cache-Aside e Estratégias de Expiração (TTL)}

% ------------------------------------------------------------------------------
\twocolframe[title={O Padrão Cache-Aside (Lazy Loading)},leftcol=7.1cm,rightcol=6.9cm]{%
  \alert{Mecânica de Funcionamento}\par
  \itemseps[topsepi=0.2mm,itemsepi=0.4mm,labelsep=1.5mm,leftmargini=2.5mm]
  \begin{itemize}
    \item \textbf{A Aplicação Coordena:} O back-end busca na memória rápida (Redis) e na base persistente (Disco/SQL).
    \item \textbf{Fluxo de Leitura:}
      \begin{enumerate}
        \item \textbf{Cache Hit:} Se a chave existir, retorna da RAM ($\approx 0.4\text{ ms}$).
        \item \textbf{Cache Miss:} Se não existir, busca na base lenta ($\approx 120\text{ ms}$), salva no Redis com \texttt{SETEX} e retorna.
      \end{enumerate}
    \item \textbf{Vantagem:} Apenas dados ativamente consultados consomem memória RAM.
  \end{itemize}
}{%
  \alert{Por que o Cache DEVE Ter TTL?}\par
  \itemseps[topsepi=0.2mm,itemsepi=0.4mm,labelsep=1.5mm,leftmargini=2.5mm]
  \begin{itemize}
    \item \textbf{Finitude de RAM:} Evita que o Redis atinja \texttt{maxmemory} e entre em evicção forçada.
    \item \textbf{Autocura:} Dados defasados (\textit{Stale Reads}) são eliminados após o tempo de vida.
    \item \textbf{Invalidação Pró-Ativa:} Ao atualizar o catálogo, deleta-se a chave (\texttt{DEL cache:catalogo}).
  \end{itemize}
}

% ------------------------------------------------------------------------------
\begin{frame}[fragile]{Implementação do Cache-Aside em Python no FastCommerce}
  \vspace{0.1cm}
  \begin{columns}[t]
    \begin{column}{0.48\textwidth}
      \footnotesize
      \alert{Lógica de Negócio na API}\par
      \vspace{0.15cm}
      \begin{itemize}\setlength{\itemsep}{1.2mm}\setlength{\parskip}{0pt}
        \item Utiliza \texttt{r.get(chave)} para consulta em memória.
        \item Grava no Redis com \texttt{r.setex(chave, 45, dados)} aplicando TTL de 45 segundos.
        \item Mensura a latência em tempo real para o painel de telemetria da interface web.
        \item Speedup comprovado: de $\approx 120\text{ ms}$ no Miss para $\approx 0.4\text{ ms}$ no Hit ($> 250\times$ mais rápido!).
      \end{itemize}
    \end{column}

    \begin{column}{0.48\textwidth}
      \alert{Código Python no FastCommerce}\par
      \vspace{0.1cm}
\begin{lstlisting}[language=Python,basicstyle=\fontsize{4.6pt}{5.6pt}\selectfont\ttfamily]
@app.route("/api/produtos")
def get_produtos():
    inicio = time.time()
    chave = "cache:catalogo_produtos"

    # 1. Verifica Cache no Redis (RAM)
    if r.exists(chave):
        produtos = json.loads(r.get(chave))
        origem = "CACHE_HIT (RAM Redis)"
    else:
        # 2. Cache Miss: Busca lenta simulada
        time.sleep(0.12)
        produtos = CATALOGO_DISCO
        # 3. Grava no Redis com TTL de 45s
        r.setex(chave, 45, json.dumps(produtos))
        origem = "CACHE_MISS (Disco/DB)"

    latencia = (time.time() - inicio) * 1000
    return jsonify({
        "origem": origem, 
        "latencia_ms": round(latencia, 2),
        "produtos": produtos
    })
\end{lstlisting}
    \end{column}
  \end{columns}
\end{frame}

\changecolours{ScoutGreen}
\section[Sessões e Carrinho com Hashes]{Controle de Sessão e Carrinho Atômico com Hashes}

% ------------------------------------------------------------------------------
\twocolframe[title={Gestão de Sessão do Usuário em Memória},leftcol=7.1cm,rightcol=6.9cm]{%
  \alert{Por que Armazenar Sessões no Redis?}\par
  \itemseps[topsepi=0.2mm,itemsepi=0.4mm,labelsep=1.5mm,leftmargini=2.5mm]
  \begin{itemize}
    \item \textbf{Escala Horizontal:} Com múltiplos servidores/containers, uma sessão local fica isolada.
    \item \textbf{Centralização em Memória:} O Redis atua como o \textit{Session Store} único compartilhado.
    \item \textbf{Expiração Deslizante (\textit{Sliding}):}
      A cada requisição, o TTL renova com \texttt{EXPIRE session:<token> 1800} (30 minutos de inatividade).
  \end{itemize}
}{%
  \alert{Hashes vs. Strings no Carrinho de Compras}\par
  \itemseps[topsepi=0.2mm,itemsepi=0.4mm,labelsep=1.5mm,leftmargini=2.5mm]
  \begin{itemize}
    \item \textbf{String (JSON Serializado):}
      Exige ler o carrinho inteiro, desserializar e salvar tudo de volta. Risco grave de \textit{Race Condition}.
    \item \textbf{Hashes do Redis:}
      A chave é o carrinho (\texttt{cart:<session>}) e os campos são os produtos (\texttt{prod\_1}, \texttt{prod\_2}).
    \item \textbf{Operação Atômica \texttt{HINCRBY}:}
      O Redis incrementa o campo diretamente na memória sem leitura prévia.
  \end{itemize}
}

% ------------------------------------------------------------------------------
\begin{frame}[fragile]{Operações Atômicas de Carrinho no FastCommerce}
  \vspace{0.1cm}
  \begin{columns}[t]
    \begin{column}{0.48\textwidth}
      \footnotesize
      \alert{Comandos Canônicos no Terminal}\par
      \vspace{0.15cm}
      \begin{itemize}\setlength{\itemsep}{1.2mm}\setlength{\parskip}{0pt}
        \item \textbf{Adicionar Item:}
          \texttt{HINCRBY cart:usr\_10 prod\_laptop 1}
        \item \textbf{Consultar Carrinho:}
          \texttt{HGETALL cart:usr\_10}
        \item \textbf{Remover Item:}
          \texttt{HDEL cart:usr\_10 prod\_mouse}
        \item \textbf{Limpar Carrinho pós-checkout:}
          \texttt{DEL cart:usr\_10}
        \item \textbf{Carrinhos Abandonados:}
          \texttt{EXPIRE cart:usr\_10 3600} elimina inativos após 1 hora.
      \end{itemize}
    \end{column}

    \begin{column}{0.48\textwidth}
      \alert{Código Python no FastCommerce}\par
      \vspace{0.1cm}
\begin{lstlisting}[language=Python,basicstyle=\fontsize{4.8pt}{5.8pt}\selectfont\ttfamily]
@app.route("/api/carrinho/adicionar", methods=["POST"])
def adicionar_ao_carrinho():
    dados = request.get_json()
    prod_id = dados.get("produto_id")
    qtd = int(dados.get("quantidade", 1))

    chave = f"cart:{session_id}"
    
    # HINCRBY: Incremento atomico no Hash
    nova_qtd = r.hincrby(chave, prod_id, qtd)
    
    # Define TTL de 1 hora no carrinho
    r.expire(chave, 3600)
    
    return jsonify({
        "status": "success",
        "produto_id": prod_id,
        "quantidade": nova_qtd
    })
\end{lstlisting}
    \end{column}
  \end{columns}
\end{frame}

\changecolours{ScoutGreen}
\section[Filas e Desacoplamento]{Filas Assíncronas de Checkout com Listas do Redis}

% ------------------------------------------------------------------------------
\twocolframe[title={O Problema do Checkout Síncrono no E-commerce},leftcol=7.1cm,rightcol=6.9cm]{%
  \alert{Por que o Checkout Síncrono Falha?}\par
  \itemseps[topsepi=0.2mm,itemsepi=0.4mm,labelsep=1.5mm,leftmargini=2.5mm]
  \begin{itemize}
    \item \textbf{Cadeia Lenta:} Cartão, antifraude, NF-e e e-mail levam de 2 a 6 segundos.
    \item \textbf{Esgotamento:} 1000 clientes simultâneos derrubam servidores por conexões abertas.
    \item \textbf{Solução com Fila:} A API valida o carrinho, empacota e joga na fila do Redis ($\approx 0.5\text{ ms}$).
  \end{itemize}
}{%
  \alert{O Padrão Produtor-Consumidor FIFO}\par
  \itemseps[topsepi=0.2mm,itemsepi=0.4mm,labelsep=1.5mm,leftmargini=2.5mm]
  \begin{itemize}
    \item \textbf{Produtor (API Web):}
      Insere na esquerda com \texttt{LPUSH fila:pedidos <dados>}.
    \item \textbf{Consumidor (Worker):}
      Retira na direita com \texttt{BRPOP fila:pedidos 5}.
    \item \textbf{Por que \texttt{BRPOP}?}\par
      Não consome CPU em loop (\textit{busy waiting}). O processo dorme no kernel até um item chegar!
  \end{itemize}
}

% ------------------------------------------------------------------------------
\begin{frame}[fragile]{Implementação do Produtor e do Consumidor}
  \vspace{0.1cm}
  \begin{columns}[t]
    \begin{column}{0.48\textwidth}
      \footnotesize
      \alert{Produtor (API Web no Checkout)}\par
      \vspace{0.1cm}
\begin{lstlisting}[language=Python,basicstyle=\fontsize{4.6pt}{5.6pt}\selectfont\ttfamily]
@app.route("/api/checkout", methods=["POST"])
def checkout():
    # 1. Le carrinho do Redis Hash
    itens = r.hgetall(f"cart:{session_id}")
    if not itens:
        return jsonify({"erro": "vazio"}), 400

    pedido = {
        "id_pedido": f"PED-{int(time.time())}",
        "session_id": session_id,
        "total": calcular_total(itens),
        "status": "ENFILEIRADO"
    }
    
    # 2. LPUSH: Enfileira pedido na lista
    pos = r.lpush("fila:pedidos", json.dumps(pedido))
    # 3. Limpa o carrinho
    r.delete(f"cart:{session_id}")
    return jsonify({"pedido": pedido, "pos": pos})
\end{lstlisting}
    \end{column}

    \begin{column}{0.48\textwidth}
      \alert{Consumidor (\texttt{worker\_pedidos.py})}\par
      \vspace{0.1cm}
\begin{lstlisting}[language=Python,basicstyle=\fontsize{4.6pt}{5.6pt}\selectfont\ttfamily]
def processar_pedidos():
    chave = "fila:pedidos"
    print("Worker escutando via BRPOP...")
    
    while True:
        # BRPOP bloqueia ate 5s sem CPU
        res = r.brpop(chave, timeout=5)
        if res is None:
            continue
            
        _, pedido_raw = res
        pedido = json.loads(pedido_raw)
        
        print(f"Processando: {pedido['id_pedido']}")
        executar_pagamento(pedido)
        
        # PREVIEW DA AULA 06 (MONGODB):
        # db.pedidos.insert_one(pedido)
\end{lstlisting}
    \end{column}
  \end{columns}
\end{frame}

\changecolours{ScoutGreen}
\section[Conclusão e Próximos Passos]{Síntese e Conexão Arquitetural com o MongoDB}

% ------------------------------------------------------------------------------
\twocolframe[title={Checklist de Competências Consolidadas},leftcol=7.1cm,rightcol=6.9cm]{%
  \alert{O Que Aprendemos e Executamos}\par
  \itemseps[topsepi=0.5mm,itemsepi=1mm,labelsep=2mm,leftmargini=3mm]
  \begin{itemize}
    \item \textbf{Padrão Cache-Aside:} Implementação com expiração automática (TTL) e mensuração do Speedup de acesso em memória ($> 250\times$).
    \item \textbf{Gerenciamento de Sessão:} Tokens efêmeros centralizados com renovação dinâmica (\textit{Sliding Expiration}).
    \item \textbf{Hashes do Redis:} Carrinho de compras atômico com \texttt{HINCRBY}, evitando condições de corrida em e-commerce.
    \item \textbf{Filas com Listas:} Desacoplamento assíncrono entre a API de checkout (\texttt{LPUSH}) e o worker de processamento (\texttt{BRPOP}).
  \end{itemize}
}{%
  \alert{Ponte com a Aula 06: O Modelo de Documentos}\par
  \itemseps[topsepi=0.5mm,itemsepi=1mm,labelsep=2mm,leftmargini=3mm]
  \begin{itemize}
    \item \textbf{Onde o Redis Termina?}\par
      O Redis é excelente para dados em trânsito e alta frequência. Porém, memória RAM é cara: não guardamos histórico fiscal de 10 anos no Redis!
    \item \textbf{O Papel do MongoDB (Aula 06):}
      \begin{enumerate}
        \item O catálogo de produtos passará a ser uma coleção flexível de documentos BSON (\texttt{produtos}).
        \item O worker que consome a fila do Redis salvará os pedidos confirmados na coleção \texttt{pedidos} no MongoDB.
      \end{enumerate}
    \item \textbf{Resultado:} Uma arquitetura poliglota completa e pronta para a AVP1!
  \end{itemize}
}

% ------------------------------------------------------------------------------
\begin{frame}{Evolução do Projeto FastCommerce rumo à AVP1}
  \vspace{0.1cm}
  \begin{center}
    \tikz[scale=0.78, transform shape]{
      % Etapa 1: Aula 05
      \node[draw=red!80!black, fill=red!10, thick, rounded corners, minimum width=3.8cm, minimum height=2.2cm, align=left] (etapa1) {
        \textbf{Aula 05 (20/10): Redis}\\
        \tiny $\bullet$ Cache-Aside com TTL\\
        \tiny $\bullet$ Sessão e Tokens em Memória\\
        \tiny $\bullet$ Carrinho de Compras em Hashes\\
        \tiny $\bullet$ Fila de Checkout com Listas
      };

      % Seta de evolucao
      \draw[->, ultra thick, ScoutNavy] (etapa1) -- node[above, font=\bfseries\tiny] {Consumo} ++(4.8cm, 0) coordinate (meio);

      % Etapa 2: Aulas 06 e 07
      \node[draw=ScoutGreen, fill=ScoutGreen!15, thick, rounded corners, minimum width=3.8cm, minimum height=2.2cm, align=left, right=0.2cm of meio] (etapa2) {
        \textbf{Aulas 06/07: MongoDB}\\
        \tiny $\bullet$ Catálogo BSON Flexível\\
        \tiny $\bullet$ Embedding de Itens/Avaliações\\
        \tiny $\bullet$ Persistência do Worker no MongoDB\\
        \tiny $\bullet$ Índices e Agregações de Vendas
      };

      % Etapa 3: AVP1
      \node[draw=ScoutNavy, fill=ScoutNavy!15, thick, rounded corners, minimum width=8.5cm, minimum height=0.9cm, below=0.6cm of etapa1, xshift=2.4cm, align=center] (avp1) {
        \textbf{Aula 10 (17/11): Avaliação Prática AVP1 em Laboratório}\\
        \tiny Sistema Web Completo com Persistência Poliglota Integrada (FastCommerce: Redis + MongoDB)
      };

      \draw[->, thick, black!70] (etapa1.south) |- (avp1.west);
      \draw[->, thick, black!70] (etapa2.south) |- (avp1.east);
    }
  \end{center}
  \vspace{0.15cm}
  \centering{\footnotesize \textbf{Roteiro Prático:} Execute o servidor web e o worker para testar o ciclo completo em laboratório!}
\end{frame}


% ------------------------------------------------------------------------------
% Encerramento
% ------------------------------------------------------------------------------
\changecolours[inverse]{ScoutGreen}
\begin{frame}[plain]
  \vspace{2cm}
  \centering
  {\Huge \textbf{Dúvidas \& Prática em Laboratório}}\par
  \vspace{0.8cm}
  {\large IFCE Campus Tauá -- ADS16 (Persistência Poliglota)}\par
  \vspace{0.3cm}
  {\normalsize FastCommerce NoSQL: Redis na Memória $\to$ MongoDB no Disco}
\end{frame}

\end{document}
